\SourceLeaf{160}mache sich aber zur Regel, nichts zu verschweigen, damit andere auch anders rechnen können, wenn sie es wünschen. Die Zahlenwerte der Ergebnisse werden, wie man auch verfährt, gleich brauchbar sein. Man läuft aber Gefahr, ihre Zuverlässigkeit zu hoch anzugeben, wenn man zu schnell mit dem Ausschließen zur Hand ist. Entscheidungen dieser Art scheinen mir eher dem \textls{Handeln im Leben} ähnlich, wo mathematische Strenge und Gewissheit selten oder nie gegeben sind und man sich damit begnügen muss, nach bestem wohlüberlegtem Urteil zu verfahren.

\begin{center}
. . . . . . . . . . . . . . . . .
\end{center}



\Needspace{6\baselineskip}
\begin{center}
\textsc{Gauss} an \textsc{Gerling}. Göttingen, 2. April 1840.
\end{center}

. . . . . . . . . . Zum wissenschaftlichen Inhalt Ihres Briefes will ich nur einige Worte zu den Bedenken sagen, die Sie sich unnötigerweise wegen der Verschiedenartigkeit von Winkel- und Längenmessungen machen.

Sind $q, q', q'', q'''$ etc. überhaupt durch Beobachtung gefundene Größen, gleichgültig, ob homogen oder heterogen, und bezeichnet man mit $dq, dq', dq''\ldots$ die Differenzen zwischen Beobachtung und Rechnung nach irgendeinem System von Elementen, so besagt das Prinzip der Methode der kleinsten Quadrate \textls{durchaus nicht}, dass
\[
dq^2+dq'^2+dq''^2+\text{etc.}
\]
ein Minimum werden soll, sondern dass
$\left(\dfrac{dq}{m}\right)^2+\left(\dfrac{dq'}{m'}\right)^2+\left(\dfrac{dq''}{m''}\right)^2+$etc.
ein Minimum werden soll, wobei $m, m', m''\ldots$ die bei den einzelnen Daten jeweils zu erwartenden mittleren Fehler bezeichnen. Nur weil in der Praxis so überwiegend oft Fälle vorkommen, in denen $m=m'=m''=m'''$ etc. gesetzt werden kann oder muss, lässt man sich dazu verleiten, die erste Aussage gelten zu lassen und zu vergessen, dass nur die zweite allgemein gültig ist.

Um dieses Prinzip anwenden zu können, muss man, wenn auch nicht genau, so doch einigermaßen näherungsweise, wenn auch nicht die Größen
